\clearpage
\VerdierPhysicalPageAnchor{0242}%
\begingroup\small\setlength{\parskip}{0pt}\setlength{\emergencystretch}{1.5em}
\setlength{\abovedisplayskip}{4pt plus 2pt minus 2pt}%
\setlength{\belowdisplayskip}{4pt plus 2pt minus 2pt}%

\VerdierSemanticAnchor{ch4d-4-5-3-b-cont}
\begin{enumerate}
\item[1)] \textit{$X$ est un objet de $D^b(\mathcal A)$.}
\item[2)] \textit{$X$ est un objet de $D^+(\mathcal A)$ et $F$ est
stationnaire à droite.}
\item[3)] \textit{$X$ est un objet de $D^-(\mathcal A)$ et $F$ est
stationnaire à gauche.}
\item[4)] \textit{Le foncteur $F$ est stationnaire.}
\end{enumerate}

\VerdierSemanticAnchor{ch4d-4-5-3-proof}
Cette proposition résulte immédiatement de (4.5.2) et de la définition
(chap.~II, 4.4.2).

\VerdierSemanticAnchor{ch4d-4-5-4}
\medskip\noindent\textbf{Théorème 4.5.4.}\quad
\textit{Soient
$F_1,F_2\colon D^*(\mathcal A)\longrightarrow\mathcal B$
($*=+,-,b,{}$ ``vide'') deux foncteurs cohomologiques (resp.\
$G_1,G_2\colon D^*(\mathcal A)\longrightarrow D^{*'}(\mathcal B)$ deux
foncteurs exacts) et $m\colon F_1\longrightarrow F_2$ (resp.\
$m\colon G_1\longrightarrow G_2$) un morphisme de foncteurs (resp.\ un
morphisme de foncteurs exacts (chap.~I, 1.6.6.4)). La sous-catégorie pleine
de $D^*(\mathcal A)$ définie par les objets $X$ de $D^*(\mathcal A)$ tels
que $m(X[p])$ soit un isomorphisme, pour tout entier $p$, est une
sous-catégorie triangulée strictement pleine et saturée de
$D^*(\mathcal A)$ (chap.~II, 2.1.6), notée $\operatorname{Is}(m)$.}

\VerdierSemanticAnchor{ch4d-4-5-4-a}
\noindent\textit{a) Si tout objet de $\mathcal A$ est un objet de
$\operatorname{Is}(m)$, alors
$D^b(\mathcal A)\subset\operatorname{Is}(m)$.}

\VerdierSemanticAnchor{ch4d-4-5-4-b}
\noindent\textit{b) Si tout objet de $\mathcal A$ est un objet de
$\operatorname{Is}(m)$ et si les foncteurs $F_1$ et $F_2$ (resp. $G_1$ et
$G_2$) sont stationnaires à droite, alors
$D^+(\mathcal A)\cap D^*(\mathcal A)\subset\operatorname{Is}(m)$.}

\VerdierSemanticAnchor{ch4d-4-5-4-c}
\noindent\textit{c) Si tout objet de $\mathcal A$ est un objet de
$\operatorname{Is}(m)$ et si les foncteurs $F_1$ et $F_2$ (resp. $G_1$ et
$G_2$) sont stationnaires à gauche, alors
$D^-(\mathcal A)\cap D^*(\mathcal A)\subset\operatorname{Is}(m)$.}

\VerdierSemanticAnchor{ch4d-4-5-4-d}
\noindent\textit{d) Si tout objet de $\mathcal A$ est un objet de
$\operatorname{Is}(m)$ et si les foncteurs $F_1$ et $F_2$ (resp. $G_1$ et
$G_2$) sont stationnaires, alors
$D^*(\mathcal A)=\operatorname{Is}(m)$.}

\VerdierSemanticAnchor{ch4d-4-5-4-e}
\noindent\textit{e) S'il existe une partie $M$ de
$\operatorname{Ob}(\mathcal A)$ telle que :}

\noindent\textit{$\alpha$) tout objet de $\mathcal A$ soit isomorphe à un
quotient d'un objet appartenant à $M$;}

\noindent\textit{$\beta$) tout objet appartenant à $M$ soit un objet de
$\operatorname{Is}(m)$;}

\noindent\textit{si de plus les foncteurs $F_1$ et $F_2$ (resp. $G_1$ et
$G_2$) sont stationnaires à gauche, alors
$D^-(\mathcal A)\cap D^*(\mathcal A)\subset\operatorname{Is}(m)$.}

\VerdierSemanticAnchor{ch4d-4-5-4-f}
\noindent\textit{f) S'il existe une partie $M$ de
$\operatorname{Ob}(\mathcal A)$ telle que :}

\noindent\textit{$\alpha$) tout objet de $\mathcal A$ soit isomorphe à un
sous-objet d'un objet appartenant à $M$;}

\noindent\textit{$\beta$) tout objet appartenant à $M$ est un objet de
$\operatorname{Is}(m)$;}

\noindent\textit{si de plus les foncteurs $F_1$ et $F_2$ (resp. $G_1$ et
$G_2$) sont stationnaires à droite, alors
$D^+(\mathcal A)\cap D^*(\mathcal A)\subset\operatorname{Is}(m)$.}

\VerdierSemanticAnchor{ch4d-4-5-4-proof-start}
Notons, tout d'abord, que la partie resp.\ du théorème (4.5.4) se déduit de
la partie non resp., en composant les foncteurs $G_1$ et $G_2$ avec le
foncteur cohomologique canonique
$D^{*'}(\mathcal B)\longrightarrow\mathcal B$. Démontrons a). Soit $X$ un
objet de
\endgroup
